% Generated draft: requires mathematical and source review.
\documentclass[11pt]{article}
\usepackage[margin=25mm]{geometry}
\usepackage{fontspec}
\setmainfont{Latin Modern Roman}
\usepackage{amsmath,amssymb,mathtools,unicode-math}
\setmathfont{Latin Modern Math}
\usepackage{enumitem,needspace,xurl,hyperref}
\setlength{\parindent}{0pt}
\setlength{\parskip}{5pt}
\setlength{\emergencystretch}{3em}
\newcommand{\N}{\mathbb{N}}
\newcommand{\Z}{\mathbb{Z}}
\newcommand{\Q}{\mathbb{Q}}
\newcommand{\R}{\mathbb{R}}
\newcommand{\C}{\mathbb{C}}
\newcommand{\F}{\mathbb{F}}
\newcommand{\A}{\mathbb{A}}
\newcommand{\PP}{\mathbb{P}}
\newcommand{\E}{\mathbb{E}}
\newcommand{\eps}{\varepsilon}
\newcommand{\dd}{\,\mathrm{d}}
\begin{document}
\section*{UnsolvedMath ID: 30003871}
\textbf{Source problem code:} catalogue entry
\subsection{Definitions and mathematical statement}
Let \(p\) be a prime, let \(E\) be a number field, let \(X\) be a smooth algebraic variety over \(E\), and let \(\sigma:E\hookrightarrow\C\) be a field embedding. Write
\[
X_\sigma=X\times_{\operatorname{Spec}E,\sigma}\operatorname{Spec}\C.
\]
Let \(L\) be a geometric \(p\)-adic étale local system on \(X\). Here a \(p\)-adic étale local system means a finite-rank lisse \(p\)-adic sheaf; the coefficient field and the additional condition denoted by ``geometric'' are not specified in the supplied record.

The source proposes a comparison of two tensor functors, one obtained through the complex Riemann--Hilbert correspondence and the other through algebraic de Rham theory. To express the intended target, provisionally denote their common source category by \(\mathcal G_p(X)\), and denote by \(\operatorname{Conn}_{\mathrm{rs}}(X_\sigma)\) the category of algebraic vector bundles on \(X_\sigma\) equipped with integrable connections having regular singularities at infinity. Integrability means that the curvature vanishes. Regularity at infinity is part of the target even when \(X\) is not proper. Write the proposed functors schematically as
\[
\mathcal R_\sigma,\mathcal D_\sigma:
\mathcal G_p(X)\longrightarrow\operatorname{Conn}_{\mathrm{rs}}(X_\sigma).
\]
These symbols name the intended constructions; their existence and precise definitions are not supplied by the record.

\textbf{Source conjecture, with unresolved construction data.} For every such \(E,X,p,\sigma\), the two proposed functors are canonically isomorphic as tensor functors:
\[
\mathcal R_\sigma\cong^{\otimes}\mathcal D_\sigma.
\]
Thus the proposed comparison must be natural in every geometric local system \(L\), compatible with tensor products and tensor units, and an isomorphism of the resulting connections. Moreover, for every such \(L\), the filtered de Rham realization is conjectured to be a complex variation of Hodge structures on \(X_\sigma\).

\textbf{Unresolved qualifications.} The record does not define \(\mathcal G_p(X)\), specify the coefficient field, or explain how \(L\) produces the complex local system required by complex Riemann--Hilbert. The embedding \(\sigma\) specifies the complex base variety but does not itself specify that coefficient comparison. Nor does the record specify the algebraic de Rham realization, its descent and comparison data, or its Hodge filtration. For a filtered connection \((\mathcal V,\nabla,F^\bullet)\), a variation of Hodge structures requires, in particular, Griffiths transversality
\[
\nabla(F^a\mathcal V)\subseteq F^{a-1}\mathcal V\otimes\Omega^1_{X_\sigma}
\qquad(a\in\Z),
\]
together with the appropriate Hodge-structure data and axioms. The weight, the convention for a complex variation, and any polarization requirement are unspecified. Consequently, the displayed assertions preserve the intended conjectural targets but do not constitute a fully determined formulation until these missing data are resolved.
\subsection{Short English statement}
For a smooth variety over a number field, the source conjectures that the complex Riemann--Hilbert and algebraic de Rham constructions attached to every geometric \(p\)-adic local system agree canonically as tensor functors. It also conjectures that the filtered de Rham realization is a complex variation of Hodge structures. The meaning of geometric, the coefficient comparisons, the realization functors, and the Hodge-theoretic conventions require clarification.
\subsection{Sources}
\begin{itemize}
\item Supplied local UnsolvedMath registry record, permanent ID \(30003871\), source problem code \texttt{catalogue entry}, title: \(p\)-Adic Riemann--Hilbert and de Rham Compatibility. Collection: Oberwolfach Reports Open Problems. The proposed year is recorded as 2018; no individual proposer is supplied.
\item Supplied report citation: \emph{Algebraische Zahlentheorie} (2019), Oberwolfach Reports, pp.\ 1785--1844, DOI \texttt{10.4171/owr/2018/30}. \url{https://doi.org/10.4171/owr/2018/30}. The report is supplied as workshop and open-problem-discussion provenance; no problem-specific page or named attribution is given.
\item Supplied literature source: \emph{Logarithmic Riemann--Hilbert correspondences}. \url{https://arxiv.org/abs/1803.05786}.
\item Supplied literature source: \emph{Oberwolfach Report 30/2018}. \url{https://ems.press/content/serial-article-files/46751}.
\item Literature assessment dated exactly \texttt{2026-08-21}: the supplied curation labels the problem open and describes the \(p\)-adic logarithmic Riemann--Hilbert comparison conjecture and its Hodge-theoretic compatibility as out of reach except in special cases, with no general proof identified. This is supplied dated curation, not an independently verified literature assessment; the background explicitly requests independent verification of time-sensitive status and cited claims.
\item The registry labels the statement \texttt{corrected\_verified} and reports that it was checked against the official Oberwolfach report and preserves its hypotheses, notation, and open target. This is a supplied source-verification claim, not verification performed for this draft.
\end{itemize}
\end{document}
